\section{Evaluation Methodology}\label{sec:method}

\paragraph{Single-margin simulator and stopping deadlines.}
The single-margin experiment models two robots on perpendicular paths. Robot A may brake and wait for robot B; an optimistic, range-dependent bias hides the conflict and additive Gaussian measurement noise perturbs the monitored margin. The bias is proportional to $\mathrm{range}^{p}$: $p<1$ reveals the conflict late, $p=1$ gives a linear reveal, and $p>1$ reveals it earlier. Recovery is rate-limited braking with capability $a_{\max}$. The predictive rule estimates the margin slope with a Theil--Sen estimator over the latest eight samples. To avoid division by a fitted slope near zero, the implementation requires a decay magnitude greater than $10^{-3}$ margin units per second; otherwise the predictive rule uses its $M\le0$ fallback.

A general point-of-no-return procedure could replay the realized future from every recorded state and determine the last state from which an admissible recovery avoids failure. We use two model-specific quantities and do not treat either as a general safety certificate. First, a continuous-time braking calculation gives the latest trigger that would stop A outside the wider crossing zone. It is conservative relative to the simulator's collision condition: with the default parameters it is $8.65$\,s ($8.6$\,s to one decimal). Second, deterministic replay on the simulator's $0.1$\,s control grid finds $8.8$\,s as the latest tested forced-trigger instant that avoids a collision in the default conflict episode. The latter is an empirical, grid-dependent boundary, not a continuous-time proof.

\paragraph{Exploratory single-margin sweeps.}
The dose-response experiment compares an uncertainty-blind arrival-gap trigger, the current-margin trigger at $M=0$, and prediction. It uses 40 paired noise seeds at each of seven severities, with the predictive horizon fixed at $T_h=2.0$\,s. The recovery-capability and degradation-shape sweeps use 25 paired seeds at each of five severities, or 125 episodes per plotted setting and method. For each setting, the fixed-pad candidates are
\[
  \rho\in\{0,0.5,1.0,1.5,2.0,2.5,3.0,4.0\}.
\]
The displayed comparator is the candidate with the lowest mean successful elapsed time among those for which every episode succeeds on that same setting and seed grid. Selection and reporting therefore use the same simulations. These figures are exploratory mechanism checks, not out-of-sample performance estimates.
